% =====================================================================
% DEPRECATED — results_spatial_diagnostics.tex
% Superseded by mmar-input/07b_spatial_diagnostics.tex, which has
% itself been moved to appendix.tex \section{Spatial Structure of Residuals}.
% Retained for git history only. Do not \input this file.
% =====================================================================
%
% These captions deliberately carry their hedges INLINE so the epistemics are
% fixed before any number lands -- do not soften them when filling values.
% =====================================================================

\subsection{Spatial validity of the representation ladder}
\label{sec:spatial-diagnostics}

The headline test (Section~\ref{sec:results}) establishes \emph{whether} the
representation ladder improves prediction under spatially-blocked validation.
This subsection asks the spatial-computing questions behind that result: are
the model's errors spatially structured, does the underlying relationship vary
across space, do the spatial features dissolve the structure they target, and
does the finding survive an alternative block geometry. All quantities here are
\textbf{descriptive}; the pre-registered fold-level Wilcoxon test remains the
sole basis for inference, and none of these diagnostics enters the geometric
compatibility index $\kappa_{\text{geom}}$.

% ---------------------------------------------------------------------
% Narrative arc (fill the bracketed values; keep the hedges):
%   1. Errors are spatially structured        -> global Moran's I + Geary's C
%   2. WHERE they cluster                      -> LISA (Fig. ref{fig:lisa})
%   3. WHY spatial structure helps             -> GWR non-stationarity (Fig. ref{fig:gwr})
%   4. Spatial features attenuate the clusters -> LISA R0->R1 (Fig. ref{fig:lisa})
%   5. Finding survives alternative blocking   -> regionalization (Table ref{tab:robustness})
% ---------------------------------------------------------------------

On the primary cell (Houston, NFIP event claims, HistGBDT), out-of-fold R0
residuals exhibit significant global spatial autocorrelation
(Moran's $I = [0.NN]$, % lisa_rollup.parquet: global_moran_I @ level=r0
$p = [0.NNN]$),
% lisa_results.json: global Moran pseudo-p (or note permutation inference)
corroborated by Geary's $C = [0.NN]$
% geary_rollup.parquet: geary_C @ r0
(Table~\ref{tab:residual-diagnostics}). Local Moran's~$I$ localizes this
structure to a contiguous High--High cluster covering $[N]\%$
% lisa_rollup.parquet: frac_HH @ r0 (as percent)
of ZCTAs, where the model systematically under-predicts
(Figure~\ref{fig:lisa}, left). Under R1's spatial features the
significant-cluster fraction falls to $[M]\%$;
% lisa_rollup.parquet: frac_HH + frac_LL @ r1
across all modelable cells, clusters attenuated in $[k]$ of $[n]$
% lisa_results.json: attenuation_summary.n_attenuated_R0_R1 / n_cells
(Clopper--Pearson 95\% CI $[[lo, hi]]$,
% lisa_results.json: attenuation_summary.clopper_pearson_95_ci
reported descriptively; $n{=}[n]$ precludes a formal test). Attenuation is
consistent with the W-matrix features absorbing the spatial error structure;
because the design is a controlled representation intervention rather than a
randomized one, we make no causal claim at the hydrologic level.

A geographically weighted regression of claims on the frozen R0 feature subset
shows the relationship is non-stationary across space: GWR improves on global
OLS ($\Delta$AICc $= [N.N]$)
% gwr_nonstationarity.json: cells[houston].gwr_delta_aicc
and local $R^2$ ranges from $[0.NN]$ to $[0.NN]$
% gwr_local_r2_houston.parquet: min/max of local_r2
(Figure~\ref{fig:gwr}). The regions of lowest local fit coincide with the LISA
clusters and, by rank, with the cells of largest R0$\rightarrow$R1 uplift
(Spearman $\rho = [0.NN]$, bootstrap 95\% CI $[[lo, hi]]$);
% gwr_nonstationarity.json: triangulation.spearman.gwr_vs_uplift
we report this as an observed triangulation across $n{=}[n]$ cells, not a
hypothesis test, and it is excluded from the Holm--Bonferroni family. GWR is fit
on the full cross-section and describes structure only; it carries no
generalization claim.

% =====================================================================
% TABLE 1 -- Residual spatial-diagnostics
% One row per (scenario, target, level). Promote the primary cell to the top.
% Source: lisa_rollup.parquet (frac_*, global_moran_I) + geary_rollup.parquet (geary_C, geary_p)
% =====================================================================
\begin{table}[t]
  \centering
  \caption{Residual spatial diagnostics by representation level. Global
    Moran's~$I$ and Geary's~$C$ on out-of-fold HistGBDT residuals (Queen
    contiguity, 999 permutations); LISA cluster fractions are FDR-controlled
    (Benjamini--Hochberg, $q{=}0.05$). Geary's~$C<1$ indicates positive
    autocorrelation. All values descriptive.}
  \label{tab:residual-diagnostics}
  \small
  \begin{tabular}{lll rr rrr}
    \toprule
    Scenario & Target & Lvl & Moran $I$ & Geary $C$ & HH & LL & Sig. \\
    \midrule
    % --- primary cell first ---
    houston & nfip & R0 & [0.NN] & [0.NN] & [0.NN] & [0.NN] & [0.NN] \\
    houston & nfip & R1 & [0.NN] & [0.NN] & [0.NN] & [0.NN] & [0.NN] \\
    houston & nfip & R2 & [0.NN] & [0.NN] & [0.NN] & [0.NN] & [0.NN] \\
    % --- remaining cells (one triple per scenario/target) ---
    % $\vdots$
    \bottomrule
  \end{tabular}
  % Column sources (per row, keyed on scenario+target+level):
  %   Moran I : lisa_rollup.parquet  -> global_moran_I
  %   Geary C : geary_rollup.parquet -> geary_C   (geary_p available for text)
  %   HH      : lisa_rollup.parquet  -> frac_HH
  %   LL      : lisa_rollup.parquet  -> frac_LL
  %   Sig.    : lisa_rollup.parquet  -> frac_significant
\end{table}

% =====================================================================
% FIGURE 5 -- LISA residual clusters (main triptych)
% Rendered by render_fig5_lisa_clusters.py --mode main
% =====================================================================
\begin{figure*}[t]
  \centering
  \includegraphics[width=\textwidth]{fig5_lisa_clusters.pdf}
  \caption{Local Moran's~$I$ on out-of-fold residuals for the primary cell
    (Houston, NFIP event claims, HistGBDT) across representation levels.
    Categories are FDR-significant (Benjamini--Hochberg, $q{=}0.05$); grey
    ZCTAs are not significant. Panel subtitles give global Moran's~$I$ and the
    significant-cluster fraction. The contiguous High--High block at R0 marks
    where the tabular model systematically under-predicts; it attenuates under
    R1's spatial features. Attenuation is consistent with, but not proof of,
    the spatial features absorbing the error structure.}
  \label{fig:lisa}
\end{figure*}

% =====================================================================
% FIGURE 6 -- GWR non-stationarity
% Rendered by render_fig6_gwr_nonstationarity.py
% NOTE: coefficient-surface panels require the gwr_params_{scenario}.parquet
%       sidecar amendment (see render_fig6 docstring). Until then the figure is
%       a single local-R^2 panel -- adjust width accordingly.
% =====================================================================
\begin{figure}[t]
  \centering
  \includegraphics[width=\columnwidth]{fig6_gwr_nonstationarity.pdf}
  \caption{Geographically weighted regression of NFIP claims on the frozen R0
    feature subset (Houston). Left: local~$R^2$, showing where the global
    relationship holds versus collapses. Remaining panels: local coefficient
    surfaces for selected covariates. GWR is fit on the full cross-section and
    describes spatial non-stationarity of the relationship; it is a structural
    diagnostic, not a measure of predictive performance, and is independent of
    $\kappa_{\text{geom}}$.}
  \label{fig:gwr}
\end{figure}

% =====================================================================
% TABLE 2 -- Regionalization robustness (Appendix R)
% Source: compute_spatial_sidecar_regionalize.py --compare output
%   header counts: n_county_cells, n_region_cells, dropped_county_only, dropped_region_only
%   rows: comparisons[] -> uplift_county_blocked, uplift_region_blocked, direction_agrees
% =====================================================================
\begin{table}[t]
  \centering
  \caption{Robustness of R0$\rightarrow$R1 uplift to block geometry. County
    blocks (registered primary) versus Skater regionalization on structural
    geography ($K{=}5$). We report sign agreement only; the county-blocked
    Wilcoxon remains the primary inference. \emph{Compared $[m]$ of $[9]$ cells}
    % compare: n_cells_compared of (n_county_cells); name dropped keys below
    --- dropped cells (regionalization degenerate or empty fold): $[\text{keys}]$.
    % compare: dropped_county_only + dropped_region_only
    Contiguous holdout raises the missing-value rate of \texttt{wlag\_nfip\_claims}
    in region-fold test ZCTAs ($[X]\%$ vs.\ $[Y]\%$ under county folds),
    % RunResult.wlag_nan_rates aggregated by split; region vs county
    a consequence of the block geometry rather than a modeling change.}
  \label{tab:robustness}
  \small
  \begin{tabular}{ll rr c}
    \toprule
    Scenario & Target & Uplift (county) & Uplift (region) & Dir.\ agrees \\
    \midrule
    houston & nfip & [+0.NN] & [+0.NN] & [yes] \\
    % $\vdots$  one row per compared cell
    \bottomrule
  \end{tabular}
  % Row sources (compare comparisons[]):
  %   Uplift (county) : uplift_county_blocked
  %   Uplift (region) : uplift_region_blocked
  %   Dir. agrees     : direction_agrees
\end{table}

We close with two robustness checks. Geary's~$C$ corroborates the Moran's~$I$
direction at every level (Table~\ref{tab:residual-diagnostics}). Under an
alternative hydrologically-coherent block definition the sign of the
R0$\rightarrow$R1 uplift agrees with county blocking in $[k]$ of $[m]$ compared
% compare: n_direction_agrees of n_cells_compared
cells (Table~\ref{tab:robustness}); where it disagrees we read this as
sensitivity to block geometry, not as a failed result.
